%%==============================================================================
%% appendix.tex — \input from report.lhs, after \appendix is called
%%==============================================================================
\section{How to use this template}
\label{app:howto}

This appendix is a quick reference for working with the template itself;
remove it from the final submission if you don't want it in the deliverable.

\subsection*{Adding a new module section}
Copy \texttt{sections/00-template.lhs}, rename it, fill in \verb|\Topic{}|,
the diagram, the code, and any \texttt{property}/\texttt{calculation}/
\texttt{unittest} blocks, then add a matching \texttt{\%include} line in
\texttt{report.lhs}.

\subsection*{Adding notation}
Math you type yourself (prose, equations, diagrams) goes in
\texttt{ap-symbols.sty}; rules controlling how literal Haskell renders go in
\texttt{ap-format.fmt}. See \autoref{tab:notation} for the symbols already
defined.

\subsection*{Figures}
Use a standard \texttt{figure} environment with \texttt{\textbackslash
includegraphics} for images, or a \texttt{tikzpicture} for diagrams drawn
directly in LaTeX (see \autoref{fig:example-tree} for an example with no
external image file needed).

\subsection*{Tables}
Use \texttt{booktabs} (\texttt{\textbackslash toprule}/\texttt{\textbackslash
midrule}/\texttt{\textbackslash bottomrule}) rather than full grid lines —
it reads as more professional and is already loaded.

\subsection*{Citations}
Add entries to \texttt{report.bib} and cite with \verb|\cite{key}|, e.g.
\verb|\cite{bird1997algebra}|.

\subsection*{Properties and proofs}
\begin{verbatim}
\begin{property}
  ... statement ...
\end{property}

\begin{calculation}
  expr
\jstep{=}{justification}
  expr'
\qedstep
\end{calculation}
\end{verbatim}

\subsection*{Unit tests (GHCi sessions)}
Write content directly inside \texttt{unittest} — no \verb|\begin{lstlisting}| wrapper needed:
\begin{verbatim}
\begin{unittest}
ghci> myFunction arg
result
\end{unittest}
\end{verbatim}
The box is automatically numbered and labelled \textbf{GHCi} with a dark
terminal style.
